\section{Reference algebraicity and rotation-word distinction}\label{app:reference}
We complete the reference-member argument from Section~\ref{sec:exact}.
For $h_0^2=1/2$, the columns defined there give the polynomial equations
\begin{equation}\label{eq:elimination}
a\cdot(b\times d)=0,\qquad
2(a\times b)_y^2-(a\times d)_y^2=0.
\end{equation}
Removing integer content and the factors $(1+u^2)^{17}$ and
$(1+u^2)^{16}$, respectively, leaves bidegrees $(22,14)$ and $(30,32)$ in
$(r,u)$. Their resultant in $r$ has degree $1094$, with factor degrees and
multiplicities
\[
(1,46),\quad(2,126),\quad(182,2),\quad(432,1).
\]
Interval evaluation in the reference box excludes the first three factors.
The stored degree-432 polynomial $f$ is checked by reconstructing the
resultant and multiplying its factors. Its irreducibility over $\mathbb Q$
is certified by squarefree, degree-preserving reductions at $7,13,37,43$:
the possible factor-degree subset sums common to all four are only $0,432$.
Thus $f$ is the minimal polynomial of the reference tilt $u$.
The first equation in~\eqref{eq:elimination} then proves $r$ algebraic:
its nonzero coefficients in $u$ have degree at most $14$, so cannot vanish
at a root of degree $432$.

An exact Gaussian-integer Rouch\'e check isolates a root $\zeta$ of $f$ in
the disk of radius $10^{-10}$ about
\[
-0.028547283375+1.138685386817\,i.
\]
After substituting the centre plus $10^{-10}w$, the linear term exceeds the
sum of all other terms by a factor greater than $100$ on $|w|=1$.
It follows that there is exactly one root in this disk. Irreducibility makes
it an algebraic conjugate of the real reference $u$.

Here is the all-length step. Regard the three oriented axes as spaced by
$120^\circ$, with each generator having positive angle $\alpha$.
Any identity between their real rotation words transfers to $\zeta$, because
the scaled-coordinate matrices are rational over $\mathbb Q(u)$ and their
denominators do not vanish there. Use the projective lifts
\[
L_a(z)=I-iz(a\cdot\boldsymbol\sigma),
\]
where $\boldsymbol\sigma$ are the Pauli matrices. Conjugation on traceless
matrices implements the complexified Rodrigues rotation; this identity is
checked symbolically. Equal rotation words thus give equal projective
transformations on $\mathbb{CP}^1$.

On its unit-sphere model, the eigenlines are $+a,-a$. The ratio of minus to
plus eigen-coordinates of a point at spherical angle $\delta$ from $+a$ has
modulus $\tan(\delta/2)$. Applying $L_a(\zeta)$ multiplies this modulus by
\[
\left|\frac{1+i\zeta}{1-i\zeta}\right|<0.067
< (2-\sqrt3)^2=\tan^2(15^\circ).
\]
Take the three disjoint closed caps of radius $30^\circ$ about the positive
axes. Each point in their union is at least $30^\circ$ from every negative
axis. Therefore each generator sends the whole union into the interior of
its own cap. The image identifies the first letter of a nonempty word.
Cancellation and induction distinguish all equal-length words.
Section~\ref{sec:signs} now proves absence of exact overlaps in the six-map
system. Hochman's theorem yields the reference dimension.
